\section{Limitations}
\label{sec:limitations}

\paragraph{Specification and attribution.}
The audit depends on task contrast, nuisance space, and representation; nuisance variation must preserve the task. Absolute eigenvalues depend on task scaling; ridge estimates also on nuisance scaling. Standardization does not establish representation invariance. Ridge inflates projected information and can conceal an unregularized collapse; a small ridge parameter controls this bias only when the nuisance block has a positive conditioning margin. Predictive Fisher weights confidence; unweighted Schur statistics describe residual feature geometry. Their relation to population efficient information depends on the score model. Centering reverses the prospective ranking in all five seeds, so its original uncentered diagnostic depends partly on slice means. Comparable CLIP and random-direction discovery results and unweighted-control intervals containing zero leave semantic nuisance projection's incremental contribution unresolved. Confidence-matched comparisons and systematic nuisance controls remain necessary; local certificates do not guarantee reliability outside audited populations.

\paragraph{Local theory and audit scope.}
Efficient information, separation strength, and local radius jointly determine classification bounds, which require regularity and quantitative Bayes-decision separation. Sample size counts the audited environment's observations. Latent-environment pooling also depends on prevalence, shared information, and discovery. Coverage is relative to the declared environment distribution. Finite audits cover included groups or frozen slice rules; small mass and weak nuisance conditioning amplify uncertainty. Bonferroni-adjusted Gaussian/Wald coverage is nominal, conditional on fitted models, evaluated families, and realized coordinate-standardization constants; uncertainty in evaluation-fitted constants is omitted. Repeated searches must account for certification-data reuse. Empirical sample-size scaling and achievable classification rates remain open.

\paragraph{Development history and empirical breadth.}
The prospective rule was frozen after pilot experiments for fresh seeds on the same benchmark. Independent confirmatory datasets and matched evaluations of comparative selection gains are needed. The probe's background attribute defines the intervention; same-type recomposition controls shared construction effects. Evidence centers on MobileNetV3-small image classifiers, supplementary CivilComments TF-IDF diagnostics (Appendix~\ref{app:predefined-diagnostics}), and synthetic constructions. Transformer, segmentation, and diffusion applications, including nuisance discovery without metadata or prior structure, require further evaluation. The unregularized training audits concern the fresh models in Appendix~\ref{app:implemented-audit}.

\paragraph{Information, performance, and intervention scope.}
Information gains need not repair predictions: prospective swap sensitivity persists, constrained slice repair selects unchanged models, and post hoc selection increases slice information while reducing worst-group accuracy. Prospective FG selection required fallbacks in seeds 22 and 24 because no candidate met ERM's overall validation accuracy minus one percentage point. Post hoc choices coincide with validation-NLL minimization. The proposed threshold-based selection policy requires prospective validation against the deployment objective. Reserves assume an informative anchor and stated information-accumulation design; a general admission procedure for nuisance relevance, task preservation, split stability, and conditioning remains to be evaluated. TTI measures local predictor sensitivity; finite interventions and retraining are empirical tests. Transfer of annotation-conditioned feature shifts to realizable label-preserving interventions and other candidate families remains unestablished.
